% ======================================================
% 第 31 天
% ======================================================
\setday{31}{定积分的换元法}

\oneproblem{
    设 $f(x)$ 是连续函数，$F(x)$ 是 $f(x)$ 的原函数，则（ \quad ）
    \begin{multicols}{2}
    \begin{enumerate}[label=(\Alph*)]
        \item 当 $f(x)$ 是奇函数时，$F(x)$ 必是偶函数
        \item 当 $f(x)$ 是偶函数时，$F(x)$ 必是奇函数
        \item 当 $f(x)$ 是周期函数时，$F(x)$ 必是周期函数
        \item 当 $f(x)$ 是单调增函数时，$F(x)$ 必是单调增函数
    \end{enumerate}
    \end{multicols}
}

\oneproblem{
    设 $f(x)$ 为连续函数，$\displaystyle I = t\int_0^{\frac{s}{t}} f(tx)\,\mathrm{d}x$，其中 $s>0, t>0$，则 $I$ 的值（ \quad ）
    \begin{multicols}{2}
    \begin{enumerate}[label=(\Alph*)]
        \item 依赖于 $s$ 和 $t$
        \item 依赖于 $s,t,x$
        \item 依赖于 $t$ 和 $x$，不依赖于 $s$
        \item 依赖于 $s$，不依赖于 $t$
    \end{enumerate}
    \end{multicols}
}

\oneproblem{
    设 $\displaystyle I_k = \int_0^{k\pi} e^{x^2}\sin x\,\mathrm{d}x\ (k=1,2,3)$，则有（ \quad ）
    \begin{multicols}{2}
    \begin{enumerate}[label=(\Alph*)]
        \item $I_1 < I_2 < I_3$
        \item $I_3 < I_2 < I_1$
        \item $I_2 < I_3 < I_1$
        \item $I_2 < I_1 < I_3$
    \end{enumerate}
    \end{multicols}
}

\oneproblem{
    设三个积分分别为$
    M = \int_{-\frac{\pi}{2}}^{\frac{\pi}{2}} \dfrac{\sin x}{1+x^2}\cos^4 x\,\mathrm{d}x,\quad
    N = \int_{-\frac{\pi}{2}}^{\frac{\pi}{2}} (\sin^3 x + \cos^4 x)\,\mathrm{d}x,\quad,P = \int_{-\frac{\pi}{2}}^{\frac{\pi}{2}} (x^2\sin^3 x - \cos^4 x)\,\mathrm{d}x,$
    则（ \quad ）
    \begin{multicols}{2}
    \begin{enumerate}[label=(\Alph*)]
        \item $N < P < M$
        \item $M < P < N$
        \item $N < M < P$
        \item $P < M < N$
    \end{enumerate}
    \end{multicols}
}

\oneproblem{
    设 $\displaystyle a_n = \dfrac{3}{2}\int_0^{\frac{n}{n+1}} x^{n-1}\sqrt{1+x^n}\,\mathrm{d}x$，则极限 $\displaystyle \lim_{n\to\infty} na_n$ 等于（ \quad ）
    \begin{multicols}{2}
    \begin{enumerate}[label=(\Alph*)]
        \item $(1+e)^{\frac{3}{2}}+1$
        \item $(1+e^{-1})^{\frac{3}{2}}-1$
        \item $(1+e^{-1})^{\frac{3}{2}}+1$
        \item $(1+e)^{\frac{3}{2}}-1$
    \end{enumerate}
    \end{multicols}
}

\oneproblem{
    设
    $
    M = \int_{-\frac{\pi}{2}}^{\frac{\pi}{2}} \dfrac{(1+x)^2}{1+x^2}\,\mathrm{d}x,\quad
    N = \int_{-\frac{\pi}{2}}^{\frac{\pi}{2}} \dfrac{1+x}{e^x}\,\mathrm{d}x,\quad
    K = \int_{-\frac{\pi}{2}}^{\frac{\pi}{2}} (1+\sqrt{\cos x})\,\mathrm{d}x,
    $
    则 $M,N,K$ 的大小关系为（ \quad ）
    \begin{multicols}{2}
    \begin{enumerate}[label=(\Alph*)]
        \item $M > N > K$
        \item $M > K > N$
        \item $K > M > N$
        \item $K > N > M$
    \end{enumerate}
    \end{multicols}
}

\oneproblem{
    求 $\displaystyle \lim_{x\to 0} \dfrac{\int_0^x (x-t)f(t)\,\mathrm{d}t}{x\int_0^x f(x-t)\,\mathrm{d}t}$，其中 $f(x)$ 连续且 $f(0)\ne 0$.
}

\oneproblem{
    设 $f(x)$ 是周期为2的连续函数：
    \begin{enumerate}[label=(\arabic*)]
        \item 证明对任意实数 $t$，有 $\displaystyle \int_t^{t+2} f(x)\,\mathrm{d}x = \int_0^2 f(x)\,\mathrm{d}x$；
        \item 证明 $\displaystyle G(x) = \int_0^x \left[2f(t) - \int_t^{t+2} f(s)\,\mathrm{d}s\right]\mathrm{d}t$ 是周期为2的周期函数.
    \end{enumerate}
}

\oneproblem{
    设 $\displaystyle f(x) = \int_1^x \dfrac{\ln t}{1+t}\,\mathrm{d}t$，其中 $x>0$，求 $\displaystyle f(x) + f\left(\dfrac{1}{x}\right)$.
}

\oneproblem{
    设 $f(x)$ 是区间 $[0,\frac{\pi}{4}]$ 上的单调、可导函数，且满足
    $
    \int_0^{f(x)} f^{-1}(t)\,\mathrm{d}t = \int_0^x t\dfrac{\cos t - \sin t}{\sin t + \cos t}\,\mathrm{d}t,
    $
    其中 $f^{-1}$ 是 $f$ 的反函数，求 $f(x)$.
}

\oneproblem{
    设 $f(x)$ 在 $(-\infty,+\infty)$ 内满足
    $
    f(x) = f(x-\pi) + \sin x,
    $
    且 $f(x)=x,\ x\in[0,\pi)$，计算 $\displaystyle I = \int_\pi^{3\pi} f(x)\,\mathrm{d}x$.
}

\oneproblem{
    设函数 $f(x)$ 可导，且 $f(0)=0$，
    $
    F(x) = \int_0^x t^{n-1}f(x^n - t^n)\,\mathrm{d}t,
    $
    求 $\displaystyle \lim_{x\to 0} \dfrac{F(x)}{x^{2n}}$.
}

\oneproblem{
    设函数 $f(x)$ 连续，$\displaystyle g(x) = \int_0^1 f(xt)\,\mathrm{d}t$，且 $\displaystyle \lim_{x\to 0} \dfrac{f(x)}{x} = A$，$A$ 为常数. 求 $g'(x)$ 并讨论 $g'(x)$ 在 $x=0$ 处的连续性.
}

\oneproblem{
    设 $f(x)$ 在 $[a,b]$ 上连续，$x\in(a,b)$，证明：
    $
    \lim_{h\to 0} \dfrac{1}{h}\int_a^x [f(t+h)-f(t)]\,\mathrm{d}t = f(x) - f(a).
    $
}

\oneproblem{
    求 $\displaystyle \lim_{x\to+\infty} \sqrt[3]{x}\int_x^{x+1} \dfrac{\sin t}{\sqrt{t+\cos t}}\,\mathrm{d}t$.
}

\oneproblem{
    设 $f(x)$ 在 $(-\infty,+\infty)$ 上具有连续导数，且 $|f(x)|\le 1,\ f'(x)>0,\ x\in(-\infty,+\infty)$，证明：对于 $0<\alpha<\beta$，成立
    $
    \lim_{n\to\infty} \int_\alpha^\beta f'\left(nx-\dfrac{1}{x}\right)\,\mathrm{d}x = 0.
    $
}

\oneproblem{
    设函数 $f(x)$ 在 $(-\infty,+\infty)$ 内连续，且
    $
    F(x) = \int_0^x (x-2t)f(t)\,\mathrm{d}t,
    $
    试证：
    \begin{enumerate}[label=(\arabic*)]
        \item 若 $f(x)$ 为偶函数，则 $F(x)$ 也是偶函数；
        \item 若 $f(x)$ 单调不增，则 $F(x)$ 单调不减.
    \end{enumerate}
}

% ======================================================
% 第 32 天
% ======================================================
\setday{32}{定积分的分部积分法}

\oneproblem{
    若
    $
    \int_{-\pi}^\pi (x - a_1\cos x - b_1\sin x)^2\,\mathrm{d}x = \min_{a,b\in\mathbb{R}} \left\{ \int_{-\pi}^\pi (x - a\cos x - b\sin x)^2\,\mathrm{d}x \right\},
    $
    则 $a_1\cos x + b_1\sin x =$（ \quad ）
    \begin{multicols}{2}
    \begin{enumerate}[label=(\Alph*)]
        \item $2\sin x$
        \item $2\cos x$
        \item $2\pi\sin x$
        \item $2\pi\cos x$
    \end{enumerate}
    \end{multicols}
}

\oneproblem{
    如下图，曲线 $C$ 的方程为 $y=f(x)$，点 $(3,2)$ 是它的一个拐点，直线 $l_1$ 与 $l_2$ 分别是曲线 $C$ 在点 $(0,0)$ 与 $(3,2)$ 处的切线，其交点为 $(2,4)$。设函数 $f(x)$ 具有三阶连续导数，计算定积分
    $
    \int_0^3 (x^2+x)f'''(x)\,\mathrm{d}x.
    $
}

\oneproblem{
    已知 $f$ 在区间 $(-1,3)$ 内有二阶连续导数，$f(0)=12,\ f(2)=2f'(2)+8$，求积分 $\displaystyle \int_0^1 xf''(2x)\,\mathrm{d}x$.
}

\oneproblem{
    计算下列定积分：
    \begin{enumerate}[label=(\arabic*),itemsep=2.5cm]
        \item $\displaystyle \int_0^1 \dfrac{x^2\arcsin x}{\sqrt{1-x^2}}\,\mathrm{d}x$.
        \item $\displaystyle \int_0^1 \dfrac{\ln(1+x)}{(2-x)^2}\,\mathrm{d}x$.
        \item $\displaystyle \int_0^{\frac{\pi}{2}} \dfrac{e^x(1+\sin x)}{1+\cos x}\,\mathrm{d}x$.
        \item $\displaystyle \int_0^{\frac{\pi}{2}} \dfrac{\cos x}{1+\tan x}\,\mathrm{d}x$.
        \item $\displaystyle \lim_{n\to\infty} \int_0^1 e^{-x}\sin nx\,\mathrm{d}x$，$n$ 为正整数.
    \end{enumerate}
}

\oneproblem{
    根据函数的定义求定积分：
    \begin{enumerate}[label=(\arabic*),itemsep=2.5cm]
        \item 已知 $\displaystyle f(x) = x\int_1^x \dfrac{\sin t^2}{t}\,\mathrm{d}t$，求 $\displaystyle \int_0^1 f(x)\,\mathrm{d}x$.
        \item 已知 $\displaystyle f(x) = \int_1^x \dfrac{\ln(t+1)}{t}\,\mathrm{d}t$，求 $\displaystyle \int_0^1 \dfrac{f(x)}{\sqrt{x}}\,\mathrm{d}x$.
    \end{enumerate}
}

\oneproblem{
    设 $f(x)$ 为连续的非负函数，当 $x\ge 0$ 时，有
    $
    \int_0^1 f(x)f(xt)\,\mathrm{d}t = \sin^2 2x,
    $
    求 $\displaystyle \int f(x)\,\mathrm{d}x$.
}

\oneproblem{
    设 $\displaystyle a_n = \int_0^1 x^n\sqrt{1-x^2}\,\mathrm{d}x\ (n=1,2,3,\cdots)$
    \begin{enumerate}[label=(\arabic*),itemsep=2.5cm]
        \item 证明：数列 $\{a_n\}$ 单调递减，且 $\displaystyle a_n = \dfrac{n-1}{n+2}a_{n-2}\ (n=2,3,\cdots)$；
        \item 求极限 $\displaystyle \lim_{n\to\infty} \dfrac{a_n}{a_{n-1}}$.
    \end{enumerate}
}

\oneproblem{
    设函数 $f(x)$ 在 $[0,\pi]$ 上连续，且
    $
    \int_0^\pi f(x)\,\mathrm{d}x = 0,\quad \int_0^\pi f(x)\cos x\,\mathrm{d}x = 0.
    $
    试证：在 $(0,\pi)$ 内至少存在两个不同的点 $\xi_1,\xi_2$，使 $f(\xi_1)=f(\xi_2)=0$.
}

\oneproblem{
    函数 $f(x)$ 在 $[0,1]$ 上具有连续导数，且
    $
    \int_0^1 f(x)\,\mathrm{d}x = \dfrac{5}{2},\quad \int_0^1 xf(x)\,\mathrm{d}x = \dfrac{3}{2},
    $
    证明：存在 $\xi\in(0,1)$，使得 $f'(\xi)=3$.
}

\oneproblem{
    设函数 $f(x)$ 的导数 $f'(x)$ 在 $[0,1]$ 上连续，$f(0)=f(1)=0$，且满足
    $
    \int_0^1 [f'(x)]^2\,\mathrm{d}x - 8\int_0^1 f(x)\,\mathrm{d}x + \dfrac{4}{3} = 0,
    $
    试求 $f(x)$ 的表达式.
}

\oneproblem{
    设 $f(x)$ 在区间 $[0,1]$ 上连续可导，且 $f(0)=f(1)=0$。求证：
    $
    \left[ \int_0^1 xf(x)\,\mathrm{d}x \right]^2 \le \dfrac{1}{45}\int_0^1 [f'(x)]^2\,\mathrm{d}x,
    $
    等号当且仅当 $f(x)=A(x-x^3)$ 时成立，其中 $A$ 是常数.
}

% ======================================================
% 第 33 天
% ======================================================
\setday{33}{具有特定结构的函数的定积分的计算}

\oneproblem{
    设函数 $f(x)$ 连续，则下列函数中，必为偶函数的是（ \quad ）
    \begin{multicols}{2}
    \begin{enumerate}[label=(\Alph*)]
        \item $\displaystyle \int_0^x f(t^2)\,\mathrm{d}t$
        \item $\displaystyle \int_0^x f^2(t)\,\mathrm{d}t$
        \item $\displaystyle \int_0^x t[f(t)-f(-t)]\,\mathrm{d}t$
        \item $\displaystyle \int_0^x t[f(t)+f(-t)]\,\mathrm{d}t$
    \end{enumerate}
    \end{multicols}
}

\oneproblem{
    设 $\displaystyle F(x) = \int_x^{x+2\pi} e^{\sin t}\sin t\,\mathrm{d}t$，则 $F(x)$ 的值（ \quad ）
    \begin{multicols}{2}
    \begin{enumerate}[label=(\Alph*)]
        \item 为正常数
        \item 为负常数
        \item 恒为零
        \item 不为常数
    \end{enumerate}
    \end{multicols}
}

\oneproblem{
    设 $\displaystyle f(x) = \begin{cases} xe^{x^2}, & -\dfrac{1}{2}\le x<\dfrac{1}{2} \\ -1, & x\ge\dfrac{1}{2} \end{cases}$，试求 $\displaystyle \int_{\frac{1}{2}}^2 f(x-1)\,\mathrm{d}x$.
}

\oneproblem{
    设 $\displaystyle f(x) = \int_x^{x+\frac{\pi}{2}} |\sin t|\,\mathrm{d}t$.
    \begin{enumerate}[label=(\arabic*),itemsep=6cm]
        \item 证明 $f(x)$ 是以 $\pi$ 为周期的周期函数；
        \item 求 $f(x)$ 的值域.
    \end{enumerate}
}

\oneproblem{
    (1) 比较 $\displaystyle \int_0^1 |\ln t|[\ln(1+t)]^n\,\mathrm{d}t$ 与 $\displaystyle \int_0^1 t^n|\ln t|\,\mathrm{d}t\ (n=1,2,\cdots)$ 的大小，说明理由.\\
    (2) 记 $\displaystyle u_n = \int_0^1 |\ln t|[\ln(1+t)]^n\,\mathrm{d}t\ (n=1,2,\cdots)$，求 $\displaystyle \lim_{n\to\infty} u_n$.
}

\oneproblem{
    设 $f(x),g(x)$ 在区间 $[-a,a]\ (a>0)$ 上连续，$g(x)$ 为偶函数，且 $f(x)$ 满足 $f(x)+f(-x)=A$（$A$ 为常数）.
    \begin{enumerate}[label=(\arabic*),itemsep=6cm]
        \item 证明：$\displaystyle \int_{-a}^a f(x)g(x)\,\mathrm{d}x = A\int_0^a g(x)\,\mathrm{d}x$；
        \item 利用(1)的结论计算定积分 $\displaystyle \int_{-\frac{\pi}{2}}^{\frac{\pi}{2}} |\sin x|\arctan e^x\,\mathrm{d}x$.
    \end{enumerate}
}

\oneproblem{
    设函数 $\displaystyle S(x) = \int_0^x |\cos t|\,\mathrm{d}t$.
    \begin{enumerate}[label=(\arabic*),itemsep=6cm]
        \item 当 $n$ 为正整数，且 $n\pi\le x<(n+1)\pi$ 时，证明：$2n\le S(x)<2(n+1)$；
        \item 求 $\displaystyle \lim_{x\to+\infty} \dfrac{S(x)}{x}$.
    \end{enumerate}
}

\oneproblem{
    设 $\displaystyle f(x) = \int_0^x \dfrac{\sin t}{\pi-t}\,\mathrm{d}t$，计算 $\displaystyle \int_0^\pi f(x)\,\mathrm{d}x$.
}

\oneproblem{
    计算下列定积分：
    \begin{enumerate}[label=(\arabic*),itemsep=6cm]
        \item $\displaystyle \int_{-2}^3 \min\left\{ \dfrac{1}{\sqrt{|x|}}, x^2, x \right\}\,\mathrm{d}x$.
        \item $\displaystyle \int_2^4 \dfrac{\sqrt{x+3}}{\sqrt{x+3}+\sqrt{9-x}}\,\mathrm{d}x$.
        \item $\displaystyle \int_0^1 e^x\left( \dfrac{1-x}{1+x^2} \right)^2\,\mathrm{d}x$.
        \item $\displaystyle \int_{-\pi}^\pi \dfrac{x\sin x\cdot\arctan e^x}{1+\cos^2 x}\,\mathrm{d}x$.
        \item $\displaystyle \int_{\frac{\pi}{4}}^{\frac{\pi}{2}} \dfrac{1+\sin x}{1+\cos x}e^x\,\mathrm{d}x$.
        \item $\displaystyle \int_0^{\frac{\pi}{2}} \dfrac{e^x(1+\sin x)}{1+\cos x}\,\mathrm{d}x$.
        \item $\displaystyle \int_{-1}^1 \dfrac{\mathrm{d}x}{(e^x+1)(1+x^2)}$.
        \item $\displaystyle \int_0^\pi x\ln(\sin x)\,\mathrm{d}x$.
        \item $\displaystyle \int_{\frac{1}{2}}^{\frac{3}{2}} \dfrac{1}{\sqrt{x-x^2}}\,\mathrm{d}x$.
        \item $\displaystyle \int_0^1 x^m(\ln x)^n\,\mathrm{d}x$（其中 $m,n$ 为正整数）.
        \item $\displaystyle \int_{e^{-2n\pi}}^1 \left| \dfrac{\mathrm{d}}{\mathrm{d}x}\cos\left( \ln\dfrac{1}{x} \right) \right|\mathrm{d}x$（其中 $n$ 为正整数）.
    \end{enumerate}
}

\oneproblem{
    设 $f(x)$ 为 $(-\infty,+\infty)$ 上连续的周期为1的周期函数，且满足 $0\le f(x)\le 1$ 与 $\displaystyle \int_0^1 f(x)\,\mathrm{d}x=1$。证明：当 $0\le x\le 13$ 时，有
    $
    \int_0^{\sqrt{x}} f(t)\,\mathrm{d}t + \int_0^{\sqrt{x+27}} f(t)\,\mathrm{d}t + \int_0^{\sqrt{13-x}} f(t)\,\mathrm{d}t \le 11,
    $
    并给出取等号的条件.
}

\oneproblem{
    设函数
    $
    f(x) = -\dfrac{1}{2}\left(1+\dfrac{1}{e}\right) + \int_{-1}^1 |x-t|e^{-t^2}\,\mathrm{d}t,
    $
    证明：在区间 $(-1,1)$ 内 $f(x)$ 有且仅有两个实根.
}

\oneproblem{
    证明积分
    $
    f(n) = \sum_{m=1}^n \int_0^m \cos\dfrac{2\pi n[x+1]}{m}\,\mathrm{d}x
    $
    的值等于 $n$ 的所有因子（包1和 $n$ 本身）之和，其中 $[x+1]$ 表示不超过 $x+1$ 的最大整数，并计算 $f(2021)$.
}

% ======================================================
% 第 34 天
% ======================================================
\setday{34}{反常积分及其敛散性的判定}

\oneproblem{
    设 $\displaystyle f(x)=\begin{cases}
    \dfrac{1}{(x-1)^{\alpha-1}}, & 1<x<e \\
    \dfrac{1}{x\ln^{\alpha+1}x}, & x\ge e
    \end{cases}$，若积分 $\displaystyle \int_1^{+\infty} f(x)\,\mathrm{d}x$ 收敛，则（ \quad ）
    \begin{multicols}{2}
    \begin{enumerate}[label=(\Alph*)]
        \item $\alpha<-2$
        \item $\alpha>2$
        \item $-2<\alpha<0$
        \item $0<\alpha<2$
    \end{enumerate}
    \end{multicols}
}

\oneproblem{
    设 $m,n$ 为正整数，则反常积分 $\displaystyle \int_0^1 \dfrac{\sqrt[m]{\ln^2(1-x)}}{\sqrt[n]{x}}\,\mathrm{d}x$ 的收敛性（ \quad ）
    \begin{multicols}{2}
    \begin{enumerate}[label=(\Alph*)]
        \item 仅与 $m$ 取值有关
        \item 仅与 $n$ 取值有关
        \item 与 $m,n$ 取值都有关
        \item 与 $m,n$ 取值都无关
    \end{enumerate}
    \end{multicols}
}

\oneproblem{
    若反常积分 $\displaystyle \int_0^{+\infty} \dfrac{1}{x^a(1+x)^b}\,\mathrm{d}x$ 收敛，则（ \quad ）
    \begin{multicols}{2}
    \begin{enumerate}[label=(\Alph*)]
        \item $a<1$ 且 $b>1$
        \item $a>1$ 且 $b>1$
        \item $a<1$ 且 $a+b>1$
        \item $a>1$ 且 $a+b>1$
    \end{enumerate}
    \end{multicols}
}

\oneproblem{
    设 $p$ 为常数，若反常积分 $\displaystyle \int_0^1 \dfrac{\ln x}{x^p(1-x)^{1-p}}\,\mathrm{d}x$ 收敛，则 $p$ 的取值范围是（ \quad ）
    \begin{multicols}{2}
    \begin{enumerate}[label=(\Alph*)]
        \item $(-1,1)$
        \item $(-1,2)$
        \item $(-\infty,1)$
        \item $(-\infty,2)$
    \end{enumerate}
    \end{multicols}
}

\oneproblem{
    计算下列反常积分：
    \begin{enumerate}[label=(\arabic*),itemsep=6cm]
        \item $\displaystyle \int_3^{+\infty} \dfrac{\mathrm{d}x}{(x-1)^4\sqrt{x^2-2x}}$.
        \item $\displaystyle \int_0^{+\infty} \dfrac{xe^{-x}}{(1+e^{-x})^2}\,\mathrm{d}x$.
        \item $\displaystyle \int_1^{+\infty} \dfrac{\mathrm{d}x}{e^{1+x}+e^{3-x}}$.
        \item $\displaystyle \int_0^{+\infty} \dfrac{\ln(1+x)}{(1+x)^2}\,\mathrm{d}x$.
        \item $\displaystyle \int_0^{+\infty} e^{-2x}|\sin x|\,\mathrm{d}x$.
        \item $\displaystyle \int_0^{+\infty} e^{-sx}x^n\,\mathrm{d}x\ (n=1,2,\cdots,s>0)$.
    \end{enumerate}
}

\oneproblem{
    已知 $\displaystyle \lim_{x\to\infty} \left(\dfrac{x-a}{x+a}\right)^x = \int_a^{+\infty} 4x^2e^{-2x}\,\mathrm{d}x$，求常数 $a$ 的值.
}

\oneproblem{
    设 $\displaystyle f(x)=\dfrac{(x+1)^2(x-1)}{x^3(x-2)}$，求 $\displaystyle I=\int_{-1}^3 \dfrac{f'(x)}{1+f^2(x)}\,\mathrm{d}x$.
}

\oneproblem{
    已知 $\displaystyle \int_0^{+\infty} \dfrac{\sin x}{x}\,\mathrm{d}x=\dfrac{\pi}{2}$，求 $\displaystyle \int_0^{+\infty} \dfrac{\sin^2 x}{x^2}\,\mathrm{d}x$.
}

\oneproblem{
    证明：$\displaystyle \lim_{x\to+\infty} \sum_{n=1}^\infty \dfrac{x}{n^2+x^2}=\dfrac{\pi}{2}$.
}

\oneproblem{
    证明广义积分 $\displaystyle \int_0^{+\infty} \dfrac{\sin x}{x}\,\mathrm{d}x$ 不是绝对收敛的.
}

\oneproblem{
    判定以下反常积分是绝对收敛还是条件收敛：
    \begin{enumerate}[label=(\arabic*),itemsep=6cm]
        \item $\displaystyle \int_0^{+\infty} e^{-ax}\sin bx\,\mathrm{d}x$ ($a,b$ 为常数且 $a>0$)；
        \item $\displaystyle \int_0^1 \dfrac{\ln x}{\sqrt{x}}\,\mathrm{d}x$；
        \item $\displaystyle \int_0^1 \dfrac{1}{\sqrt{x}}\sin\dfrac{1}{x}\,\mathrm{d}x$.
    \end{enumerate}
}

\oneproblem{
    设 $f(x)$ 在 $[0,+\infty)$ 上连续，并且无穷积分 $\displaystyle \int_0^{+\infty} f(x)\,\mathrm{d}x$ 收敛. 求 $\displaystyle \lim_{y\to+\infty} \dfrac{1}{y}\int_0^y xf(x)\,\mathrm{d}x$.
}

\oneproblem{
    讨论 $\displaystyle \int_0^{+\infty} \dfrac{x}{\cos^2 x + x^\alpha\sin^2 x}\,\mathrm{d}x$ 的敛散性，其中 $\alpha$ 是一个实常数.
}

\oneproblem{
    设 $f\in C^1[0,+\infty)$，$f(0)>0,\ f'(x)\ge 0,\ \forall x\in[0,+\infty)$. 已知 $\displaystyle \int_0^{+\infty} \dfrac{1}{f(x)+f'(x)}\,\mathrm{d}x<+\infty$，证明：$\displaystyle \int_0^{+\infty} \dfrac{1}{f(x)}\,\mathrm{d}x<+\infty$.
}

\oneproblem{
    证明积分 $\displaystyle \int_0^{+\infty} \dfrac{\sqrt{x}\cos x}{x+3}\,\mathrm{d}x$ 条件收敛.
}

\oneproblem{
    讨论积分 $\displaystyle \int_1^{+\infty} \dfrac{\sin x}{x^p}\,\mathrm{d}x\ (p>0)$ 的敛散性，如果收敛，是条件收敛还是绝对收敛？
}

\oneproblem{
    证明积分 $\displaystyle \int_1^{+\infty} \dfrac{\sin x}{x}\left(1+e^{-x}\right)\,\mathrm{d}x$ 收敛.
}

% ======================================================
% 第 35 天
% ======================================================
\setday{35}{积分等式命题的证明}

\oneproblem{
    (1) 若 $f(x)$ 在 $[-a,a]$ 上连续，则
    $
    \int_{-a}^a f(x)\,\mathrm{d}x = \int_0^a [f(-x)+f(x)]\,\mathrm{d}x.
    $
    (2) 计算 $\displaystyle \int_{-\frac{\pi}{4}}^{\frac{\pi}{4}} \dfrac{1}{1+\sin x}\,\mathrm{d}x$.
}

\oneproblem{
    (1) 设函数 $f(x)$ 在 $[0,\pi]$ 上连续，则有
    $
    \int_0^\pi xf(\sin x)\,\mathrm{d}x = \dfrac{\pi}{2}\int_0^\pi f(\sin x)\,\mathrm{d}x.
    $\\
    (2) 计算 $\displaystyle \int_0^\pi \dfrac{x\sin x}{1+\sin x}\,\mathrm{d}x$.
}

\oneproblem{
    (1) 设函数 $f(x)$ 连续，则
    $
    \int_a^b f(x)\,\mathrm{d}x = \int_a^b f(a+b-x)\,\mathrm{d}x.
    $\\
    (2) 计算 $\displaystyle \int_{\frac{\pi}{6}}^{\frac{\pi}{3}} \dfrac{\cos^2 x}{x(\pi-2x)}\,\mathrm{d}x$.
}

\oneproblem{
    设 $\lambda\in\mathbb{R}$，求证：
    $
    \int_0^{\frac{\pi}{2}} \dfrac{1}{1+(\tan x)^\lambda}\,\mathrm{d}x = \int_0^{\frac{\pi}{2}} \dfrac{1}{1+(\cot x)^\lambda}\,\mathrm{d}x = \dfrac{\pi}{4}.
    $
}

\oneproblem{
    设 $f(x)$ 为 $[0,1]$ 上的连续函数，证明：
    $
    \int_0^{\frac{\pi}{2}} f(|\cos x|)\,\mathrm{d}x = \dfrac{1}{4}\int_0^{2\pi} f(|\cos x|)\,\mathrm{d}x.
    $
}

\oneproblem{
    设 $f(x)$ 为连续函数，且满足 $f(x)+f(y)=f(xy),\ x>0,y>0$. 证明：
    $
    \int_0^1 \dfrac{f(1+x)}{1+x^2}\,\mathrm{d}x = \dfrac{\pi}{8}f(2).
    $
}

\oneproblem{
    设函数 $f(x)$ 在区间 $[a,b]$ 上连续且严格递增，证明：
    $
    \int_a^b f(x)\,\mathrm{d}x = bf(b)-af(a)-\int_{f(a)}^{f(b)} f^{-1}(x)\,\mathrm{d}x.
    $
}

\oneproblem{
    若 $f(x)$ 在 $[a,b]$ 上连续，$g(x)$ 在 $[a,b]$ 上单调且可导，则至少存在一点 $\xi\in[a,b]$，使得
    $
    \int_a^b f(x)g(x)\,\mathrm{d}x = g(a)\int_a^\xi f(x)\,\mathrm{d}x + g(b)\int_\xi^b f(x)\,\mathrm{d}x.
    $
}

\oneproblem{
    设函数 $f(x)$ 在区间 $[0,1]$ 上连续且 $\displaystyle \int_0^1 f(x)\,\mathrm{d}x\ne 0$，证明：在区间 $[0,1]$ 上存在三个不同的点 $x_1,x_2,x_3$，使得
    $
    \dfrac{\pi}{8}\int_0^1 f(x)\,\mathrm{d}x = \left[\dfrac{1}{1+x_1^2}\int_0^{x_1} f(t)\,\mathrm{d}t + f(x_1)\arctan x_1\right]x_3 = \left[\dfrac{1}{1+x_2^2}\int_0^{x_2} f(t)\,\mathrm{d}t + f(x_2)\arctan x_2\right](1-x_3).
    $
}

\oneproblem{
    若函数 $f(x)$ 在 $[a,b]$ 上一阶连续可导，且 $f(a)=0$。证明：至少存在一点 $\xi\in[a,b]$，使得
    $
    f'(\xi) = \dfrac{2}{(b-a)^2}\int_a^b f(x)\,\mathrm{d}x.
    $
}

\oneproblem{
    设函数 $f(x)$ 在 $[a,b]$ 上具有二阶连续导数。求证：至少存在一点 $\xi\in[a,b]$，使得
    $
    \int_a^b f(x)\,\mathrm{d}x = f\left(\dfrac{a+b}{2}\right)(b-a) + \dfrac{(b-a)^3}{24}f''(\xi).
    $
}

\oneproblem{
    设 $f(x)$ 在 $[a,b]$ 上连续，证明：
    $
    2\int_a^b f(x)\left[\int_x^b f(t)\,\mathrm{d}t\right]\mathrm{d}x = \left[\int_a^b f(x)\,\mathrm{d}x\right]^2.
    $
}

\oneproblem{
    设 $f(x)$ 在区间 $[-a,a]\ (a>0)$ 上具有二阶连续导数，$f(0)=0$.
    \begin{enumerate}[label=(\arabic*)]
        \item 写出 $f(x)$ 的带拉格朗日余项的一阶麦克劳林公式；
        \item 证明：在 $[-a,a]$ 上至少存在一点 $\eta$，使
        $
        a^3f''(\eta)=3\int_{-a}^a f(x)\,\mathrm{d}x.
        $
    \end{enumerate}
}

\oneproblem{
    设 $f(x)$ 在 $[a,b]$ 上有二阶导数，且 $f''(x)$ 在 $[a,b]$ 上黎曼可积。证明：
    $
    f(x)=f(a)+f'(a)(x-a)+\int_a^x (x-t)f''(t)\,\mathrm{d}t,\ \forall x\in[a,b].
    $
}

\oneproblem{
    (1) 证明：函数方程 $x^3-3x=t$ 存在三个在闭区间 $[-2,2]$ 上连续，在开区间 $(-2,2)$ 内连续可微的解 $x=\varphi_1(t),x=\varphi_2(t),x=\varphi_3(t)$ 满足：
    $
    \varphi_1(-t)=-\varphi_3(t),\ \varphi_2(-t)=-\varphi_2(t),\ |t|\le 2.
    $
    (2) 若 $f$ 是 $[-2,2]$ 上的连续偶函数，证明：
    $
    \int_1^2 f(x^3-3x)\,\mathrm{d}x = \int_0^1 f(x^3-3x)\,\mathrm{d}x.
    $
}

% ======================================================
% 第 36 天
% ======================================================
\setday{36}{积分不等式命题的证明}

\oneproblem{
    设 $f(x),g(x)$ 在 $[a,b]$ 上连续，且满足
    $
    \int_a^x f(t)\,\mathrm{d}t \ge \int_a^x g(t)\,\mathrm{d}t,\ x\in[a,b),\quad \int_a^b f(t)\,\mathrm{d}t = \int_a^b g(t)\,\mathrm{d}t.
    $
    证明：$\displaystyle \int_a^b xf(x)\,\mathrm{d}x \le \int_a^b xg(x)\,\mathrm{d}x$.
}

\oneproblem{
    设 $f(x),g(x)$ 在 $[0,1]$ 上的导数连续，且 $f(0)=0,\ f'(x)\ge 0,\ g'(x)\ge 0$。证明：对任何 $a\in[0,1]$，有
    $
    \int_0^a g(x)f'(x)\,\mathrm{d}x + \int_a^1 f(x)g'(x)\,\mathrm{d}x \ge f(a)g(1).
    $
}

\oneproblem{
    设函数 $f(x)$ 在 $(-\infty,+\infty)$ 有二阶连续导数，证明：$f''(x)\ge 0$ 的充要条件为对不同实数 $a,b$，
    $
    f\left(\dfrac{a+b}{2}\right) \le \dfrac{1}{b-a}\int_a^b f(x)\,\mathrm{d}x.
    $
}

\oneproblem{
    设 $f(x)$ 在 $[0,1]$ 上可导，$f(0)=0$，且当 $x\in(0,1)$，$0<f'(x)<1$。试证：当 $a\in(0,1)$ 时，有
    $
    \left(\int_0^a f(x)\,\mathrm{d}x\right)^2 > \int_0^a f^3(x)\,\mathrm{d}x.
    $
}

\oneproblem{
    证明：对于连续函数 $f(x)>0$，有
    $
    \ln\int_0^1 f(x)\,\mathrm{d}x \ge \int_0^1 \ln f(x)\,\mathrm{d}x.
    $
}

\oneproblem{
    设 $f(x)$ 在区间 $[0,1]$ 上连续，且 $1\le f(x)\le 3$。证明：
    $
    1 \le \int_0^1 f(x)\,\mathrm{d}x \int_0^1 \dfrac{1}{f(x)}\,\mathrm{d}x \le \dfrac{4}{3}.
    $
}

\oneproblem{
    若 $f$ 在 $[a,b]$ 上连续且为下凸函数，则 $\forall x_1,x_2\in[a,b],\ x_2>x_1$，有
    $
    f\left(\dfrac{x_1+x_2}{2}\right) \le \dfrac{1}{x_2-x_1}\int_{x_1}^{x_2} f(x)\,\mathrm{d}x \le \dfrac{f(x_1)+f(x_2)}{2}.
    $
}

\oneproblem{
    设 $f(x)$ 在 $[a,b]$ 上连续可导，且 $f(a)=f(b)=0$。证明：
    $
    \int_a^b f^2(x)\,\mathrm{d}x \le \dfrac{(b-a)^2}{4}\int_a^b [f'(x)]^2\,\mathrm{d}x.
    $
}

\oneproblem{
    设函数 $f(x)$ 在 $[0,1]$ 上连续，且 $\displaystyle \int_0^1 f(x)\,\mathrm{d}x=1$。证明：
    $
    \int_0^1 (1+x^2)f^2(x)\,\mathrm{d}x \ge \dfrac{4}{\pi}.
    $
}

\oneproblem{
    设函数 $f(x)$ 在 $[a,b]$ 上连续可导，且 $f(a)=0$。证明：
    $
    \max_{x\in[a,b]} |f(x)| \le \sqrt{b-a}\cdot\sqrt{\int_a^b [f'(x)]^2\,\mathrm{d}x}.
    $
}

\oneproblem{
    证明：$\displaystyle \int_0^1 \dfrac{\cos x}{\sqrt{1-x^2}}\,\mathrm{d}x > \int_0^1 \dfrac{\sin x}{\sqrt{1-x^2}}\,\mathrm{d}x$.
}

\oneproblem{
    设 $|f(x)|\le\pi,\ f'(x)\ge m>0\ (a\le x\le b)$，证明：
    $
    \left|\int_a^b \sin f(x)\,\mathrm{d}x\right| \le \dfrac{2}{m}.
    $
}

\oneproblem{
    设 $f(x)$ 是 $[0,+\infty)$ 上非负可导函数，$f(0)=0,\ f'(x)\le\dfrac{1}{2}$。假设 $\displaystyle \int_0^{+\infty} f(x)\,\mathrm{d}x$ 收敛。求证：对于任意 $\alpha>1$，$\displaystyle \int_0^{+\infty} f^\alpha(x)\,\mathrm{d}x$ 也收敛，并且
    $
    \int_0^{+\infty} f^\alpha(x)\,\mathrm{d}x \le \left(\int_0^{+\infty} f(x)\,\mathrm{d}x\right)^\beta,\ \beta=\dfrac{\alpha+1}{2}.
    $
}

\oneproblem{
    设 $0<f(x)<1$，无穷积分 $\displaystyle \int_0^{+\infty} f(x)\,\mathrm{d}x$ 和 $\displaystyle \int_0^{+\infty} xf(x)\,\mathrm{d}x$ 都收敛。求证：
    \[
    \int_0^{+\infty} xf(x)\,\mathrm{d}x > \dfrac{1}{2}\left(\int_0^{+\infty} f(x)\,\mathrm{d}x\right)^2.
    \]
}

\oneproblem{
    设 $f:\mathbb{R}\to(0,+\infty)$ 是一可微函数，且对所有 $x,y\in\mathbb{R}$，有 $|f'(x)-f'(y)|\le|x-y|^\alpha$，其中 $\alpha\in(0,1]$ 是常数。求证：对所有 $x\in\mathbb{R}$，有 $|f'(x)|^{\frac{\alpha+1}{\alpha}} < \dfrac{\alpha+1}{\alpha}f(x)$.
}

\oneproblem{
    设 $f(x),g(x)$ 是 $[0,1]$ 区间上的单调递增函数，满足 $0\le f(x),g(x)\le 1,\ \displaystyle \int_0^1 f(x)\,\mathrm{d}x=\int_0^1 g(x)\,\mathrm{d}x$。求证：
    $
    \int_0^1 |f(x)-g(x)|\,\mathrm{d}x \le \dfrac{1}{2}.
    $
}

% ======================================================
% 第 37 天
% ======================================================
\setday{37}{定积分在几何学上的应用}

\oneproblem{
    设 $f(x),g(x)$ 在区间 $[a,b]$ 上连续，且 $g(x)<f(x)<m$（$m$ 为常数），则曲线 $y=g(x),y=f(x),x=a$ 及 $x=b\ (a<b)$ 所围成图形绕直线 $y=m$ 旋转而成的旋转体体积为（ \quad ）
    \begin{multicols}{2}
    \begin{enumerate}[label=(\Alph*)]
        \item $\displaystyle \int_a^b \pi[2m-f(x)+g(x)][f(x)-g(x)]\,\mathrm{d}x$
        \item $\displaystyle \int_a^b \pi[2m-f(x)-g(x)][f(x)-g(x)]\,\mathrm{d}x$
        \item $\displaystyle \int_a^b \pi[m-f(x)+g(x)][f(x)-g(x)]\,\mathrm{d}x$
        \item $\displaystyle \int_a^b \pi[m-f(x)-g(x)][f(x)-g(x)]\,\mathrm{d}x$
    \end{enumerate}
    \end{multicols}
}

\oneproblem{
    过曲线 $y=\sqrt[3]{x}\ (x\ge 0)$ 上的点 $A$ 作切线，使得该切线与曲线及 $x$ 轴所围成的平面图形的面积为 $\dfrac{3}{4}$，求点 $A$ 的坐标.
}

\oneproblem{
    设 $\Gamma$ 为抛物线，$P$ 是与焦点位于抛物线同侧的一点。过 $P$ 的直线 $L$ 与 $\Gamma$ 围成的有界区域的面积记作 $A(L)$。证明：$A(L)$ 取最小值当且仅当 $P$ 恰为 $L$ 被 $\Gamma$ 所截出的线段的中点.
}

\oneproblem{
    设抛物线 $y=ax^2+bx+2\ln c$ 过原点，当 $0\le x\le 1$ 时，$y\ge 0$，又已知该抛物线与 $x$ 轴及直线 $x=1$ 所围图形的面积为 $\dfrac{1}{3}$。试确定 $a,b,c$，使此图形绕 $x$ 轴旋转一周而成的旋转体的体积 $V$ 最小.
}

\oneproblem{
    设函数 $f(x)$ 在 $(0,+\infty)$ 上有定义，且满足
    $
    2f(x)+x^2f\left(\dfrac{1}{x}\right)=\dfrac{x^2+2x}{\sqrt{1+x^2}}.
    $
    \begin{enumerate}[label=(\arabic*)]
        \item 求 $f(x)$；
        \item 求曲线 $y=f(x),\ y=\dfrac{1}{2},\ y=\dfrac{\sqrt{3}}{2}$ 及 $y$ 轴围成的图形绕 $x$ 轴旋转一周所得旋转体的体积.
    \end{enumerate}
}

\oneproblem{
    设位于曲线 $y=\dfrac{1}{\sqrt{x(1+\ln^2x)}}\ (e\le x<+\infty)$ 下方，$x$ 轴上方的无界区域为 $G$，试计算 $G$ 绕 $x$ 轴旋转一周所得空间区域的体积.
}

\oneproblem{
    计算曲线 $L_1:y=\dfrac{1}{3}x^3+2x\ (0\le x\le 1)$ 绕直线 $L_2:y=\dfrac{4}{3}x$ 旋转所生成的旋转曲面的面积.
}

\oneproblem{
    已知曲线 $L:\begin{cases} x=f(t) \\ y=\cos t \end{cases}\ (0\le t<\dfrac{\pi}{2})$，其中函数 $f(t)$ 具有连续导数，且 $f(0)=0,\ f'(t)>0\ (0\le t<\dfrac{\pi}{2})$。若曲线 $L$ 的切线与 $x$ 轴的交点到切点的距离恒为 $1$，求函数 $f(t)$ 的表达式，并求以曲线 $L$ 与 $x$ 轴与 $y$ 轴为边界的区域的面积.
}

\oneproblem{
    设函数 $f(x)=\dfrac{x}{1+x},\ x\in[0,1]$，定义函数列：
    $
    f_1(x)=f(x),\ f_2(x)=f(f_1(x)),\ \cdots,\ f_n(x)=f(f_{n-1}(x)),\ \cdots
    $
    记 $S_n$ 是由曲线 $y=f_n(x)$、直线 $x=1$ 及 $x$ 轴所围平面图形的面积，求极限 $\displaystyle \lim_{n\to\infty} nS_n$.
}

\oneproblem{
    求曲线 $y=e^{-x}\sin x\ (x\ge 0)$ 与 $x$ 轴所围图形的面积.
}

\oneproblem{
    设函数 $f(x)$ 在区间 $[a,b]$ 上连续，且在 $(a,b)$ 内有 $f'(x)>0$。证明：在 $(a,b)$ 内存在唯一的 $\xi$，使曲线 $y=f(x)$ 与两直线 $y=f(\xi),x=a$ 所围成的平面图形的面积 $S_1$ 是曲线 $y=f(x)$ 与两直线 $y=f(\xi),x=b$ 所围成的平面图形的面积 $S_2$ 的 $3$ 倍.
}

\oneproblem{
    设曲线方程为 $y=e^{-x}\ (x\ge 0)$.
    \begin{enumerate}[label=(\arabic*)]
        \item 把曲线 $y=e^{-x}$、$x$ 轴、$y$ 轴和直线 $x=\xi\ (\xi>0)$ 所围平面图形绕 $x$ 轴旋转一周，得一旋转体，求此旋转体体积 $V(\xi)$；并求满足 $\displaystyle V(a)=\dfrac{1}{2}\lim_{\xi\to+\infty}V(\xi)$ 的 $a$.
        \item 在此曲线上找一点，使过该点的切线与两个坐标轴所夹平面图形的面积最大，并求出该面积.
    \end{enumerate}
}

\oneproblem{
    已知一抛物线通过 $x$ 轴上的两点 $A(1,0),B(3,0)$.
    \begin{enumerate}[label=(\arabic*)]
        \item 求证：两坐标轴与该抛物线所围图形的面积等于 $x$ 轴与该抛物线所围图形的面积；
        \item 计算上述两个平面图形绕 $x$ 轴旋转一周所产生的两个旋转体体积之比.
    \end{enumerate}
}

\oneproblem{
    设 $y=f(x)$ 是区间 $[0,1]$ 上的任一非负连续函数.
    \begin{enumerate}[label=(\arabic*)]
        \item 试证：存在 $x_0\in(0,1)$，使得在区间 $[0,x_0]$ 上以 $f(x_0)$ 为高的矩形面积等于在区间 $[x_0,1]$ 上以 $y=f(x)$ 为曲边的曲边梯形面积.
        \item 又设 $f(x)$ 在区间 $(0,1)$ 内可导，且 $f'(x)>-\dfrac{2f(x)}{x}$，证明(1)中的 $x_0$ 是唯一的.
    \end{enumerate}
}

\oneproblem{
    设 $\rho=\rho(x)$ 是抛物线 $y=\sqrt{x}$ 上任一点 $M(x,y)\ (x\ge 1)$ 处的曲率半径，$s=s(x)$ 是该抛物线上介于点 $A(1,1)$ 与 $M$ 之间的弧长，计算 $\displaystyle 3\rho\dfrac{\mathrm{d}^2\rho}{\mathrm{d}s^2}-\left(\dfrac{\mathrm{d}\rho}{\mathrm{d}s}\right)^2$ 的值.
}

\oneproblem{
    设位于第一象限的曲线 $y=f(x)$ 过点 $\left(\dfrac{\sqrt{2}}{2},\dfrac{1}{2}\right)$，其上任一点 $P(x,y)$ 处的法线与 $y$ 轴的交点为 $Q$，且线段 $PQ$ 被 $x$ 轴平分.
    \begin{enumerate}[label=(\arabic*)]
        \item 求曲线 $y=f(x)$ 的方程；
        \item 已知曲线 $y=\sin x$ 在 $[0,\pi]$ 上的弧长为 $l$，试用 $l$ 表示曲线 $y=f(x)$ 的弧长 $s$.
    \end{enumerate}
}

\oneproblem{
    曲线 $y=\dfrac{e^x+e^{-x}}{2}$ 与直线 $x=0,x=t\ (t>0)$ 及 $y=0$ 围成一曲边梯形。该曲边梯形绕 $x$ 轴旋转一周得一旋转体，其体积为 $V(t)$，侧面积为 $S(t)$，在 $x=t$ 处的底面积为 $F(t)$.
    \begin{enumerate}[label=(\arabic*)]
        \item 求 $\dfrac{S(t)}{V(t)}$ 的值；
        \item 计算极限 $\displaystyle \lim_{t\to+\infty}\dfrac{S(t)}{F(t)}$.
    \end{enumerate}
}

\oneproblem{
    如图，$C_1$ 和 $C_2$ 分别是 $y=\dfrac{1}{2}(1+e^x)$ 和 $y=e^x$ 的图形，过点 $(0,1)$ 的曲线 $C_3$ 是一单调增函数的图形。过 $C_2$ 上任一点 $M(x,y)$ 分别作垂直于 $x$ 轴和 $y$ 轴的直线 $l_x$ 和 $l_y$。记 $C_1,C_2$ 与 $l_x$ 所围图形的面积为 $S_1(x)$；$C_2,C_3$ 与 $l_y$ 所围图形的面积为 $S_2(y)$。如果总有 $S_1(x)=S_2(y)$，求曲线 $C_3$ 的方程 $x=\phi(y)$.
}

\oneproblem{
    已知曲线 $L:\begin{cases} x=f(t) \\ y=\cos t \end{cases}\ (0\le t<\dfrac{\pi}{2})$，其中函数 $f(t)$ 具有连续导数，且 $f(0)=0,\ f'(t)>0\ (0\le t<\dfrac{\pi}{2})$。若曲线 $L$ 的切线与 $x$ 轴的交点到切点的距离恒为 $1$，求函数 $f(t)$ 的表达式，并求以曲线 $L$ 与 $x$ 轴与 $y$ 轴为边界的区域的面积.
}

% ======================================================
% 第 38 天
% ======================================================
\setday{38}{定积分在物理上的应用}

\oneproblem{
    甲乙两人赛跑，计时开始时，甲在乙前方10（单位:m）处。图中，实线表示甲的速度曲线 $v=v_1(t)$（单位:m/s）虚线表示乙的速度曲线 $v=v_2(t)$，三块阴影部分面积的数值依次为10，20，3，计时开始后乙追上甲的时刻记为 $t_0$（单位:s），则（ \quad ）
    \begin{multicols}{2}
    \begin{enumerate}[label=(\Alph*)]
        \item $t_0=10$
        \item $15<t_0<20$
        \item $t_0=25$
        \item $t_0>25$
    \end{enumerate}
    \end{multicols}
}

\oneproblem{
    如图，$x$ 轴上有一线密度为常数 $\mu$，长度为 $l$ 的细杆，若质量为 $m$ 的质点到杆右端的距离为 $a$，已知引力系数为 $k$，则质点和细杆之间引力的大小为（ \quad ）
    \begin{multicols}{2}
    \begin{enumerate}[label=(\Alph*)]
        \item $\displaystyle \int_{-l}^0 \dfrac{km\mu\,\mathrm{d}x}{(a-x)^2}$
        \item $\displaystyle \int_0^l \dfrac{km\mu\,\mathrm{d}x}{(a-x)^2}$
        \item $\displaystyle 2\int_{-l/2}^0 \dfrac{km\mu\,\mathrm{d}x}{(a+x)^2}$
        \item $\displaystyle 2\int_0^{l/2} \dfrac{km\mu\,\mathrm{d}x}{(a+x)^2}$
    \end{enumerate}
    \end{multicols}
}

\oneproblem{
    斜边长为 $2a$ 的等腰直角三角形平板铅直地沉没在水中，且斜边与水面相齐，记重力加速度为 $g$，水的密度为 $\rho$，试计算三角形平板的一侧受到的水压力.
}

\oneproblem{
    一根长度为 $1$ 的细棒位于 $x$ 轴的区间 $[0,1]$ 上，若其线密度 $\rho(x)=-x^2+2x+1$，试计算该细棒的质心坐标 $\bar{x}$.
}

\oneproblem{
    为清除井底的污泥，用缆绳将抓斗放入井底，抓起污泥后提出井口（见图），已知井深30m，抓斗自重400N，缆绳每米50N，抓斗抓起的污泥重2000N，提升速度为3m/s，在提升的过程中，污泥以20N/s的速率从缝隙中漏掉，现将抓起污泥的抓斗提升至井口，问克服重力需做多少焦耳的功？（说明： $1\,\text{N}\times1\,\text{m}=1\,\text{J}$；m, N, s, J分别表示米，牛顿，秒，焦耳；\ding{173} 抓斗的高度位于井口上方的缆绳长度忽略不计）.
}

\oneproblem{
    某闸门的形状与大小如下图所示，其中直线 $l$ 为对称轴，闸门的上部为矩形 $ABCD$，下部由二次抛物线与线段 $AB$ 所围成。当水面与闸门的上端相平时，欲使闸门矩形部分承受的水压力与闸门下部承受的水压力之比为 $5:4$，闸门矩形部分的高 $h$ 应为多少米？
}

\oneproblem{
    某建筑工程打地基时，需用汽锤将桩打进土层。汽锤每次击打，都将克服土层对桩的阻力而作功。设土层对桩的阻力的大小与桩被打进地下的深度成正比(比例系数为 $k,k>0$)。汽锤第一次击打将桩打进地下 $a$ m。根据设计方案，要求汽锤每次击打桩时所作的功与前一次击打时所作的功之比为常数 $r(0<r<1)$。问：
    \begin{enumerate}[label=(\arabic*)]
        \item 汽锤击打桩3次后，可将桩打进地下多深？
        \item 若击打次数不限，汽锤至多能将桩打进地下多深？(注:m表示长度单位米.)
    \end{enumerate}
}

\oneproblem{
    一个高为 $l$ 的柱体形贮油罐，底面是长轴为 $2a$，短轴为 $2b$ 的椭圆。现将贮油罐以长轴平行于水平面平放，当油罐中油面高度为 $\dfrac{3}{2}b$ 时，计算油的质量。(长度单位为m，质量单位为kg，油的密度为常数 $\rho\,\text{kg/m}^3$)
}

\oneproblem{
    一容器的内侧是由曲线 $C$ 绕 $y$ 轴旋转一周而成的曲面，其中曲线 $C$ 由 $x^2+y^2=2y\ (y\ge\dfrac{1}{2})$ 与 $x^2+y^2=1\ (y\le\dfrac{1}{2})$ 连接而成.
    \begin{enumerate}[label=(\arabic*)]
        \item 求容器的容积；
        \item 若将容器内盛满的水从容器顶部全部抽出，至少需要做多少功？(长度单位:m，重力加速度为 $g\,\text{m/s}^2$，水的密度为 $10^3\,\text{kg/m}^3$).
    \end{enumerate}
}

% ======================================================
% 第 39 天
% ======================================================
\setday{39}{微分方程的基本概念}

\oneproblem{
    已知 $y=\dfrac{x}{\ln x}$ 是微分方程 $y'=\dfrac{y}{x}+\varphi\left(\dfrac{x}{y}\right)$ 的解，试求 $\varphi\left(\dfrac{x}{y}\right)$ 的表达式.
}

\oneproblem{
    若 $y=(1+x^2)^2-\sqrt{1+x^2},\ y=(1+x^2)^2+\sqrt{1+x^2}$ 是微分方程 $y'+p(x)y=q(x)$ 的两个解，试求 $q(x)$ 的表达式.
}

\oneproblem{
    判断函数 $y=C_1e^{2x}+C_2e^{-2x}$ 是否为微分方程 $y''-4y=0$ 的通解.
}

\oneproblem{
    写出以下列函数为通解的微分方程，其中 $C,C_1,C_2$ 为任意常数：
    \begin{enumerate}[label=(\arabic*)]
        \item $y=x+Cx^2$；
        \item $y=x^2+C_1\ln x+C_2$.
    \end{enumerate}
}

\oneproblem{
    求以 $a,b,c$ 为自由参数的曲线族 $y=\dfrac{ax+b}{x+c}$ 为通解的微分方程.
}

\oneproblem{
    求平面上的圆族所对应的方程所满足的微分方程.
}

\oneproblem{
    设一抛物线族的对称轴为 $y$ 轴，且具有公共焦点 $(0,0)$，求以描述此曲线的方程为通解的微分方程.
}

\oneproblem{
    （导弹跟踪飞机问题）设在初始时刻 $t=0$ 时导弹位于坐标原点 $(0,0)$，飞机位于点 $(a,b)$，飞机沿着平行于 $x$ 轴方向以常速 $v_0$ 飞行。导弹在时刻 $t$ 的位置为点 $(x,y)$，且速度为常值 $v_1\ (v_1>v_0)$。导弹在飞行过程中，按照制导系统始终指向飞机。试建立导弹飞行的微分方程模型.
}

% ======================================================
% 第 40 天
% ======================================================
\setday{40}{可分离变量的微分方程与齐次方程}

\oneproblem{
    若连续函数 $f(x)$ 满足关系式 $f(x)=\int_0^{2x} f\left(\dfrac{t}{2}\right)\,\mathrm{d}t+\ln2$，试求 $f(x)$ 的表达式.
}

\oneproblem{
    设函数 $f(x)$ 满足 $f(x+\Delta x)-f(x)=2xf(x)\Delta x+o(\Delta x)$，且 $f(0)=2$，试求 $f(1)$ 的值.
}

\oneproblem{
    求在 $[0,+\infty)$ 上的可微函数 $f(x)$，使 $f(x)=e^{-u(x)}$，其中 $u=\int_0^x f(t)\,\mathrm{d}t$.
}

\oneproblem{
    求下列微分方程的通解：
    \begin{enumerate}[label=(\arabic*)]
        \item $(3x^2+2xy-y^2)\mathrm{d}x+(x^2-2xy)\mathrm{d}y=0$.
        \item $\dfrac{\mathrm{d}y}{\mathrm{d}x}=(y-x)^2+2$.
        \item $\dfrac{\mathrm{d}y}{\mathrm{d}x}=\dfrac{2x^3+3xy^2-7x}{3x^2y+2y^3-8y}$.
        \item $(2x+y-4)\mathrm{d}x+(x+y-1)\mathrm{d}y=0$.
    \end{enumerate}
}

\oneproblem{
    求下列微分方程初值问题（特解）：
    \begin{enumerate}[label=(\arabic*)]
        \item $\begin{cases} x^2y'+xy=y^2, \\ y|_{x=1}=1. \end{cases}$
        \item $\begin{cases} (y+\sqrt{x^2+y^2})\mathrm{d}x-x\mathrm{d}y=0\ (x>0), \\ y|_{x=1}=0. \end{cases}$
        \item $\begin{cases} (x+1)\dfrac{\mathrm{d}y}{\mathrm{d}x}+1=2e^{-y}, \\ y(0)=0. \end{cases}$
        \item $\begin{cases} \dfrac{\mathrm{d}y}{\mathrm{d}x}=|y+1|+|y-1|, \\ y(0)=0. \end{cases}$
    \end{enumerate}
}

\oneproblem{
    设曲线 $L$ 的极坐标方程为 $r=r(\theta)$，$M(r,\theta)$ 为 $L$ 上任一点，$M_0(2,0)$ 为 $L$ 上一定点，若极径 $OM_0,OM$ 与曲线 $L$ 所围成的曲边扇形面积值等于 $L$ 上 $M_0,M$ 两点间弧长值的一半，求曲线 $L$ 的直角坐标方程.
}

\oneproblem{
    设 $f(x)$ 在 $[1,+\infty)$ 上连续，若由曲线 $y=f(x)$，直线 $x=1,x=t\ (t>1)$ 与 $x$ 轴所围成的平面图形绕 $x$ 轴旋转一周所成的旋转体体积为
    $
    V(t)=\dfrac{\pi}{3}\left[t^2f(t)-f(1)\right],
    $
    试求 $y=f(x)$ 所满足的微分方程，并求该微分方程满足条件 $y|_{x=2}=\dfrac{2}{9}$ 的解.
}

\oneproblem{
    设 $f(x)$ 是区间 $[0,+\infty)$ 上具有连续导数的单调增加函数，且 $f(0)=1$。对任意的 $t\in[0,+\infty)$，直线 $x=0,x=t$，曲线 $y=f(x)$ 以及 $x$ 轴所围成的曲边梯形绕 $x$ 轴旋转一周生成一旋转体，若该旋转体的侧面面积在数值上等于其体积的2倍，求函数 $f(x)$ 的表达式.
}

\oneproblem{
    设函数 $y=y(x)$ 由参数方程 $\begin{cases} x=x(t) \\ y=\int_0^{t^2} \ln(1+u)\,\mathrm{d}u \end{cases}$ 确定，其中 $x(t)$ 是初值问题 $\begin{cases} \dfrac{\mathrm{d}x}{\mathrm{d}t}-2te^{-x}=0 \\ x|_{t=0}=0 \end{cases}$ 的解，求 $\dfrac{\mathrm{d}^2y}{\mathrm{d}x^2}$.
}

\oneproblem{
    已知函数 $y=y(x)$ 满足微分方程 $x^2+y^2y'=1-y'$，且 $y(2)=0$，求 $y(x)$ 的极大值与极小值.
}

\oneproblem{
    设函数 $f(x)$ 在定义域 $I$ 上的导数大于零，若对任意的 $x_0\in I$，曲线 $y=f(x)$ 在点 $(x_0,f(x_0))$ 处切线与直线 $x=x_0$ 及 $x$ 轴所围成的区域的面积恒为 $4$，且 $f(0)=2$，求 $f(x)$ 的表达式.
}

\oneproblem{
    设 $y(x)$ 是区间 $\left(0,\dfrac{3}{2}\right)$ 内的可导函数，且 $y(1)=0$，点 $P$ 是曲线 $L:y=y(x)$ 上的任意一点，$L$ 在点 $P$ 处的切线与 $y$ 轴相交于点 $(0,Y_P)$，法线与 $x$ 轴相交于点 $(X_P,0)$，若 $X_P=Y_P$，求 $L$ 上点的坐标 $(x,y)$ 满足的方程.
}